\section{Ton ETF maison}\label{sec:portefeuille}
% Chapitre 6 (tache 10, passage chapitre 06). Figures: fig_aufteilungen (deux
% figures: capital et revenu), fig_anzahl_titel, fig_sektoren, fig_laender,
% fig_rendite_wachstum (nuage de points ecrit ici, hors \linienfigur/\balkenfigur),
% fig_div_beispiele; tableau data/tab_portefeuille.tex. Aucun chiffre tape: macros de
% data/kennzahlen.tex (kapitel_6 et _macros_redaction_kapitel_6 dans
% scripts/rechnung_retraite.py). Les seuils de la methode sont lus dans
% scripts/portefeuille_exemple.py (_regles_selection), jamais recopies.
% Decision (tache 10): fig_aufteilungen reste calculee a \SparquoteZwanzig % d'epargne,
% taux commun aux cinq strategies; le texte et les legendes le disent, et les ages au
% taux de reference viennent des macros du chapitre 5 (tab:age-depart-strategies).

\noindent\textbf{Portefeuille construit par r\`egles pour illustrer la m\'ethode, pas une
recommandation d'achat.} Ce chapitre r\'epond \`a la question~: quelle part de ton capital confier \`a
ton ETF maison, combien de titres y mettre et comment les choisir~? Le
chapitre~\ref{sec:accumulation} s'est conclu sur un constat d\'efavorable~: \`a
\SparquoteZwanzig~\% d'\'epargne, aucune strat\'egie de dividendes n'atteint l'objectif, et \`a
\ReferenzSparquote~\% la r\'epartition de base n'y parvient qu'\`a l'\^age l\'egal. Ce chapitre ouvre
la bo\^ite de l'ETF maison, dont le rendement de \RenditeDivMaison~\% et la r\'epartition par pays ont
nourri tous les calculs pr\'ec\'edents. La section~\ref{sec:porte-repartitions} compare les trois
r\'epartitions entre ETF maison et ETF monde, la section~\ref{sec:porte-diversification} mesure ce
que la diversification apporte selon le nombre de titres, et la section~\ref{sec:porte-methode}
d\'ecrit la m\'ethode de s\'election et ses rel\^achements. Les titres retenus, leurs secteurs et leurs
pays sont pr\'esent\'es \`a la section~\ref{sec:porte-titres}, leur rendement face \`a la croissance
de leur dividende \`a la section~\ref{sec:porte-rendement}, et l'historique r\'eel de deux dividendes
\`a la section~\ref{sec:porte-exemples}. La section~\ref{sec:porte-biais} d\'eclare les biais de
l'exercice.

\subsection{Quelle part confier \`a l'ETF maison~?}\label{sec:porte-repartitions}

Le plan compare trois r\'epartitions de l'\'epargne~: moiti\'e dans l'ETF maison et moiti\'e dans un
ETF monde distribuant, la r\'epartition de base \ReferenzStrategie{}, et tout dans l'ETF maison. Les
deux strat\'egies de r\'ef\'erence en ETF monde, avec retrait de \TauxRetrait~\% ou en part
distribuante, compl\`etent la comparaison.

\annahme{M\^emes hypoth\`eses que le chapitre~\ref{sec:accumulation}~: rendement r\'eel total de
\RenditeReal~\% par an pour les deux ETF (\RenditeReelleConvention), dont \RenditeDivMaison~\% de
dividende pour l'ETF maison et \RenditeDivEtf~\% pour l'ETF monde. L'ETF maison n'est donc pas suppos\'e
battre le march\'e~: son dividende cro\^it de \CroissanceModeleMaison~\% par an en termes r\'eels dans
le mod\`ele, bien moins que la croissance pass\'ee de ses titres (section~\ref{sec:porte-rendement}).
Taux d'\'epargne commun de \SparquoteZwanzig~\% du salaire net pour les cinq strat\'egies, et non le
taux de \ReferenzSparquote~\% du sc\'enario de base, capital de d\'epart de \KapitalDepartMitte~€,
comparaison \`a la derni\`ere ann\'ee de l'horizon, \AnneeHorizonFin. \HypotheseUmschichtung.}

Ce taux commun de \SparquoteZwanzig~\% est un choix. Au taux du sc\'enario de base, les trois
r\'epartitions atteignent toutes l'objectif et le tableau~\ref{tab:age-depart-strategies} donne
d\'ej\`a leurs \^ages~: \AlterMaisonFuenfzigDreissig~ans pour la moiti\'e en maison,
\AlterYannBasis~ans pour la r\'epartition \ReferenzStrategie{} et \AlterMaisonHundertDreissig~ans pour
tout en maison. \`A \SparquoteZwanzig~\%, aucune ne l'atteint, et comparer le capital et le revenu
r\'eunis \`a la fin de l'horizon montre ce que l'\^age seul ne dit pas. La
figure~\ref{fig:capital-selon-repartition} donne ce capital.

% Ordre des etiquettes = ordre des lignes de data/fig_aufteilungen.csv (MaisonFuenfzig,
% MaisonSiebzig, MaisonHundert, EtfReferenz, EtfUmschichtung).
\balkenfigur{fig_aufteilungen}[ylabel={Capital en \AnneeHorizonFin{} (millions d'euros de 2026)},
  y filter/.expression={y/1e6}, ymin=0, enlarge y limits={upper, value=0.2}, bar width=22pt,
  xticklabels={{Moiti\'e en maison},{70~\% en maison},{Tout en maison},{ETF, retrait \TauxRetrait~\%},
    {ETF distribuant}},
  nodes near coords, every node near coord/.append style={font=\scriptsize,
    /pgf/number format/fixed, /pgf/number format/fixed zerofill, /pgf/number format/precision=2}]%
  {strategie}{kapital_xlv}%
  {Capital r\'euni en \AnneeHorizonFin{} selon la strat\'egie, \`a \SparquoteZwanzig~\% d'\'epargne
  (et non au taux de \ReferenzSparquote~\% du sc\'enario de base), capital de d\'epart de
  \KapitalDepartMitte~€, euros de 2026.\label{fig:capital-selon-repartition}}

Ce graphique montre que plus la part en maison est grande, plus le capital final est bas~: la
moiti\'e en maison r\'eunit \KapitalAufteilungMaisonFuenfzig~€, la r\'epartition \ReferenzStrategie{}
\KapitalHorizonZwanzig~€ et tout en maison \KapitalAufteilungMaisonHundert~€~; la moiti\'e en maison
d\'epasse ainsi tout en maison de \KapitalAufteilungEcartProzent~\%. C'est le co\^ut fiscal du
chapitre~\ref{sec:fiscalite}~: les dividendes d'actions d\'etenues en direct sont impos\'es chaque
ann\'ee sans exon\'eration partielle, et l'imp\^ot pay\'e pendant l'accumulation ne capitalise plus.
L'ETF distribuant affiche le capital le plus \'elev\'e, \KapitalAufteilungEtfUmschichtung~€, mais la
section~\ref{sec:fisc-kumuliert} a signal\'e que le mod\`ele sous-estime son imp\^ot~; cet \'ecart ne
doit pas \^etre lu comme un avantage. La figure~\ref{fig:revenu-selon-repartition} traduit ces
capitaux en revenu net.

\balkenfigur{fig_aufteilungen}[ylabel={Revenu net en \AnneeHorizonFin{} (euros de 2026 par mois)},
  ymin=0, enlarge y limits={upper, value=0.15}, bar width=22pt, scaled y ticks=false,
  extra y ticks={\ZielMonatNombre}, extra y tick labels={},
  extra y tick style={grid=major, major grid style={black, dashed}},
  xticklabels={{Moiti\'e en maison},{70~\% en maison},{Tout en maison},{ETF, retrait \TauxRetrait~\%},
    {ETF distribuant}},
  nodes near coords, every node near coord/.append style={font=\scriptsize, anchor=north,
    /pgf/number format/fixed, /pgf/number format/precision=0}]%
  {strategie}{netto_monat_xlv}%
  {Revenu net mensuel soutenable en \AnneeHorizonFin{} selon la strat\'egie, apr\`es imp\^ot et
  assurance maladie, \`a \SparquoteZwanzig~\% d'\'epargne (euros de 2026). La ligne en tirets marque
  l'objectif de \ZielMonat~€.\label{fig:revenu-selon-repartition}}

Ce graphique montre qu'\`a \SparquoteZwanzig~\% d'\'epargne, les trois r\'epartitions restent sous la
ligne de l'objectif en \AnneeHorizonFin~: \RevenuAufteilungMaisonFuenfzig~€ par mois pour la moiti\'e
en maison (\RevenuAufteilungMaisonFuenfzigPartCible~\% de la cible), \RevenuHorizonZwanzig~€ pour la
r\'epartition \ReferenzStrategie{} (\RevenuHorizonZwanzigPartCible~\%) et \RevenuAufteilungMaisonHundert~€
pour tout en maison (\RevenuAufteilungMaisonHundertPartCible~\%). L'ordre s'inverse par rapport au
capital~: tout en maison verse \RevenuAufteilungEcartProzent~\% de plus que la moiti\'e, parce que
l'ETF maison distribue \RenditeDivMaison~\% contre \RenditeDivEtf~\% pour l'ETF monde, et ce surcro\^it
de dividende l'emporte sur l'imp\^ot plus lourd. La r\'ef\'erence en ETF avec retrait de
\TauxRetrait~\% change d'\'echelle~: elle a atteint l'objectif \`a \AlterEtfReferenzZwanzig~ans et
permettrait en \AnneeHorizonFin{} de retirer \RevenuAufteilungEtfReferenz~€ nets par mois,
\RevenuAufteilungEtfReferenzMultiple~fois la cible, avec un capital du m\^eme ordre que celui des
r\'epartitions en dividendes. \bk{La part confi\'ee \`a l'ETF maison ne d\'eplace le revenu que de
\RevenuAufteilungEcartProzent~\% au plus~; c'est la r\`egle de revenu, dividendes seuls ou retrait de
\TauxRetrait~\%, qui d\'ecide de l'essentiel.}

\subsection{Combien de titres faut-il~?}\label{sec:porte-diversification}

Un ETF maison n'a d'int\'er\^et que s'il r\'epartit le risque. La
figure~\ref{fig:risque-nombre-titres} mesure, sur les titres du portefeuille-exemple eux-m\^emes, la
fluctuation d'un portefeuille selon le nombre de titres qu'il contient.

\annahme{Rendements mensuels des cours, hors dividendes, de \DonneesDebut{} \`a \DonneesFin{},
t\'el\'echarg\'es une seule fois pour la s\'election. Pour chaque nombre de titres $n$,
\TiragesDiversification{} portefeuilles \'equipond\'er\'es de $n$ titres sont tir\'es au hasard parmi les
\NombreTitres{} titres retenus (graine fixe), et l'\'ecart-type annualis\'e de leur rendement est
moyenn\'e. Aucune valeur n'est reprise de la litt\'erature.}

\linienfigur{fig_anzahl_titel}[xlabel={Nombre de titres dans le portefeuille},
  ylabel={\'Ecart-type annualis\'e du rendement (\%)}, y filter/.expression={y*1e2},
  legend pos=north east, legend cell align=left, legend style={fill=white}]%
  {n}{risiko}%
  {{Portefeuille \'equipond\'er\'e}}%
  {Fluctuation annuelle d'un portefeuille \'equipond\'er\'e selon le nombre de titres tir\'es parmi
  les \NombreTitres{} titres du portefeuille-exemple (rendements mensuels des cours,
  \DonneesDebut--\DonneesFin).\label{fig:risque-nombre-titres}}

Ce graphique montre que l'\'ecart-type passe de \RisqueUnTitre~\% pour un titre seul \`a
\RisqueTousTitres~\% pour les \NombreTitres{} titres r\'eunis, soit \RisqueReductionProzent~\% de
fluctuation en moins. L'essentiel du gain vient des premiers titres~: \`a \TitresNeufDixiemes~titres,
neuf dixi\`emes de cette baisse sont d\'ej\`a acquis. Le reste, environ \RisqueTousTitres~\%, ne
dispara\^it pas en ajoutant des titres semblables~: c'est le risque commun au march\'e, que la
diversification ne supprime pas. Cette mesure a deux limites. Elle porte sur des titres d\'ej\`a
choisis pour leur r\'egularit\'e, donc plus homog\`enes que le march\'e, et elle ne couvre que
\RegleHistoriqueAns{} ann\'ees, qui commencent apr\`es la crise financi\`ere de 2008. Un ETF monde
r\'epartit ce m\^eme risque sur un bien plus grand nombre de titres, et un d\'etracteur de la
strat\'egie des dividendes lui reproche justement ses d\'efauts de
diversification\QP{finanznachrichten_irrweg}.

\subsection{La m\'ethode de s\'election}\label{sec:porte-methode}

La s\'election applique des r\`egles chiffr\'ees fix\'ees avant de regarder le r\'esultat, sans choix
titre par titre. L'univers de d\'epart r\'eunit \UniversTitres{} titres~: \UniversAristocrates{}
aristocrates du dividende am\'ericains, \UniversDax{} valeurs du DAX et \UniversEurostoxx{} valeurs de
l'EURO STOXX~50 qui ne sont pas d\'ej\`a dans le DAX\footnote{Listes des membres lues sur Wikip\'edia
(source secondaire)~: \href{https://en.wikipedia.org/wiki/S\%26P_500_Dividend_Aristocrats}{\IndiceActionsUs{}
Dividend Aristocrats}, \href{https://en.wikipedia.org/wiki/DAX}{DAX} et
\href{https://en.wikipedia.org/wiki/EURO_STOXX_50}{EURO STOXX~50}, consult\'ees le
\DateDonneesPortefeuille.}. Cours et dividendes mensuels couvrent \RegleHistoriqueAns~ans et
proviennent d'un seul t\'el\'echargement group\'e\footnote{Yahoo Finance, cours et dividendes mensuels
et donn\'ees de soci\'et\'e (secteur, taux de distribution, flux de tr\'esorerie libre), lus par la
biblioth\`eque yfinance le \DateDonneesPortefeuille{},
\href{https://finance.yahoo.com}{finance.yahoo.com} (source secondaire).}.

\annahme{Rendement du dividende sur les douze derniers mois au cours du jour, taux de distribution,
flux de tr\'esorerie libre et secteur tels que publi\'es le \DateDonneesPortefeuille~; un titre sans
donn\'ee de taux de distribution ou de flux de tr\'esorerie est \'ecart\'e. Une baisse du dividende est
une baisse annuelle de plus de \RegleBaisseSeuil~\%.}

Un titre est retenu s'il a vers\'e au moins \RegleDividendesMin~dividendes sur la p\'eriode, si son
rendement se situe entre \RegleRendementMin{} et \RegleRendementMax~\%, si son dividende n'a pas
baiss\'e sur les \RegleAnneesCroissance{} derni\`eres ann\'ees et s'il a cr\^u d'au moins
\RegleCroissanceMin~\% par an sur la m\^eme p\'eriode. Il faut encore que la soci\'et\'e distribue au plus
\ReglePayoutMax~\% de son b\'en\'efice et que son flux de tr\'esorerie libre couvre au moins
\RegleCouvertureFcf~fois les dividendes vers\'es. Le rendement minimal \'ecarte les titres qui ne
paient presque rien, le maximum ceux dont le rendement \'elev\'e signale souvent un dividende menac\'e
(section~\ref{sec:capital-rendement}), et les deux derniers crit\`eres v\'erifient que le dividende
est pay\'e par l'activit\'e plut\^ot que par la dette. Les titres qui passent sont class\'es par un
score, moyenne de leur rang en rendement net pour un r\'esident allemand et de leur rang en
croissance du dividende, puis retenus dans cet ordre, \`a raison de \RegleSecteurMax~titres au plus
par secteur.

La m\'ethode visait \NombreTitresVise~titres~; elle en retient \NombreTitres. L'entonnoir explique
pourquoi. Sur les \UniversTitres{} titres de d\'epart, \UniversSansHistorique{} n'ont pas assez de
dividendes dans les donn\'ees, \UniversFiltreRendement{} passent le filtre du rendement et de
l'absence de baisse, et \NombreTitres{} seulement franchissent les crit\`eres de croissance, de
distribution, de couverture et de secteur. Deux rel\^achements ont \'et\'e essay\'es, un seul \`a la
fois~: une croissance d'au moins \RegleCroissanceRelache~\% au lieu de \RegleCroissanceMin~\%, puis
un taux de distribution d'au plus \ReglePayoutRelache~\% au lieu de \ReglePayoutMax~\%. Ils donnent
respectivement \PalierCroissanceTitres{} et \PalierPayoutTitres{} titres, pas un de plus. Des
\TitresEcartesFin{} titres \'ecart\'es au dernier filtre, \TitresEcartesCroissanceOk{} avaient pourtant
une croissance suffisante et \TitresEcartesSansRecul{} n'avaient pas \RegleAnneesCroissance~ans de
cotation~: ce sont donc surtout le taux de distribution, la couverture par le flux de tr\'esorerie ou
le plafond par secteur qui les ont exclus. Les donn\'ees conserv\'ees ne disent pas lequel de ces
trois crit\`eres a jou\'e pour chaque titre~; on sait seulement que \SecteursAuPlafond{} secteurs,
\SecteursAuPlafondNoms, atteignent le plafond de \RegleSecteurMax~titres.

\bk{Il manque donc \NombreTitresManquants{} titres, et ce document les laisse manquer plut\^ot que
d'assouplir les r\`egles qui font la m\'ethode.} Admettre un dividende d\'ej\`a baiss\'e ou un
rendement hors de la fourchette changerait la nature du portefeuille pour atteindre un nombre rond.
La figure~\ref{fig:risque-nombre-titres} montre d'ailleurs que \NombreTitres{} titres donnent
presque toute la diversification que ces titres peuvent offrir. Le score, lui, a \'et\'e calcul\'e
avec la retenue am\'ericaine sans formulaire W-8BEN, alors que ce document suppose ce formulaire
d\'epos\'e (section~\ref{sec:fisc-quellensteuer})~; comme moins de \NombreTitresVise{} titres ont pass\'e
les filtres, il n'a d\'epartag\'e que les titres des secteurs au plafond.

\subsection{Les titres retenus}\label{sec:porte-titres}

Le tableau~\ref{tab:titres-portefeuille-exemple} donne les \NombreTitres{} titres retenus, dans
l'ordre du score. Le rendement net y est celui d'un r\'esident allemand apr\`es retenue \`a la source
et imp\^ot allemand, avant assurance maladie et sans forfait d'\'epargnant.

\annahme{\HypotheseUsFormulaire\QH{quellensteuer}~; retenue fran\c{c}aise au taux standard, comme au
chapitre~\ref{sec:fiscalite}. La croissance est le taux annuel moyen du dividende sur
\RegleAnneesCroissance~ans, en monnaie de cotation et en termes nominaux.}

\begin{table}[htbp]
  \centering\small
  \caption{Les \NombreTitres{} titres du portefeuille-exemple, class\'es par score d\'ecroissant~:
  rendement brut sur douze mois, rendement net pour un r\'esident allemand, croissance annuelle du
  dividende sur \RegleAnneesCroissance~ans et taux de distribution (donn\'ees du
  \DateDonneesPortefeuille).}\label{tab:titres-portefeuille-exemple}
  \begin{tabularx}{\linewidth}{>{\bfseries}l l X *{4}{r}}
    \toprule
    \rowcolor[gray]{0.9}
    \textbf{Titre} & \textbf{Pays} & \textbf{Secteur} & \textbf{Rendement} & \textbf{Net}
      & \textbf{Croissance} & \textbf{Distribution} \\
    \midrule
    \tabellenkoerper{data/tab_portefeuille}
    \bottomrule
  \end{tabularx}
\end{table}

Ce tableau montre des rendements bruts de \RenditeTitreMin~\% (\RenditeTitreMinTicker) \`a
\RenditeTitreMax~\% (\RenditeTitreMaxTicker), avec une m\'ediane de \RenditeDivMaison~\%, celle qui
alimente tout le plan. Apr\`es retenue et imp\^ot, la m\'ediane nette tombe \`a \RenditeNetteMediane~\%~:
avec le formulaire W-8BEN, le pays compte peu, puisqu'il reste \NettoUs~€ pour cent euros bruts
am\'ericains, \NettoDe~€ pour des allemands et \NettoFr~€ pour des fran\c{c}ais
(chapitre~\ref{sec:fiscalite}). Le taux de distribution m\'edian est de \PayoutMedian~\%~: la plupart
des soci\'et\'es gardent une marge pour maintenir leur dividende lors d'une mauvaise ann\'ee.

La figure~\ref{fig:secteurs-portefeuille-exemple} r\'epartit les titres par secteur, chaque titre
pesant \PoidsParTitre~\% du portefeuille.

\balkenfigur{fig_sektoren}[ylabel={Part des titres (\%)}, y filter/.expression={y*1e2}, ymin=0,
  enlarge y limits={upper, value=0.2}, bar width=18pt,
  nodes near coords, every node near coord/.append style={font=\scriptsize,
    /pgf/number format/fixed, /pgf/number format/fixed zerofill, /pgf/number format/precision=1}]%
  {libelle}{anteil}%
  {R\'epartition sectorielle des \NombreTitres{} titres du portefeuille-exemple, en part du nombre
  de titres (\'equipond\'eration).\label{fig:secteurs-portefeuille-exemple}}

Ce graphique montre \SecteursNombre{} secteurs, dont \SecteursAuPlafondNoms{} au plafond de
\RegleSecteurMax~titres, soit \SecteurPartMax~\% chacun~; le plus petit p\`ese \SecteurPartMin~\%. Le
plafond par secteur a donc jou\'e son r\^ole~: sans lui, les secteurs les plus r\'eguliers auraient
pris davantage de place. Aucun titre des services aux collectivit\'es, de l'immobilier ou de la
communication n'est retenu, bien que les services aux collectivit\'es soient un secteur classique des
portefeuilles de dividendes. La figure~\ref{fig:pays-portefeuille-exemple} donne la r\'epartition par
pays.

% Les rotations TikZ s'additionnent: rotate=-35 annule le rotate=35 par defaut de
% \balkenfigur (etiquettes de pays courtes, laissees horizontales).
\balkenfigur{fig_laender}[ylabel={Part des titres (\%)}, y filter/.expression={y*1e2}, ymin=0,
  enlarge y limits={upper, value=0.2}, bar width=30pt, x tick label style={rotate=-35, anchor=north},
  nodes near coords, every node near coord/.append style={font=\scriptsize,
    /pgf/number format/fixed, /pgf/number format/fixed zerofill, /pgf/number format/precision=1}]%
  {libelle}{anteil}%
  {R\'epartition par pays de domiciliation des \NombreTitres{} titres du portefeuille-exemple, en part
  du nombre de titres.\label{fig:pays-portefeuille-exemple}}

Ce graphique montre que \PartUs~\% des titres sont am\'ericains, \PartDe~\% allemands et
\PartFr~\% fran\c{c}ais~: \PortefeuillePays{} pays seulement, alors que l'univers en comptait
\UniversPays. C'est cette r\'epartition qui fixe les retenues \`a la source de la cascade du
chapitre~\ref{sec:depart}. Un portefeuille aussi concentr\'e d\'epend des r\`egles fiscales et de la
conjoncture de trois pays, ce qu'un ETF monde dilue.

\subsection{Rendement ou croissance du dividende}\label{sec:porte-rendement}

Un dividende peut \^etre \'elev\'e aujourd'hui ou cro\^itre vite, rarement les deux \`a la fois. La
figure~\ref{fig:rendement-contre-croissance} place chaque titre retenu selon ces deux grandeurs.

\begin{figure}[htbp]
  \centering
  \begin{tikzpicture}
    \begin{axis}[retraite, width=0.95\linewidth,
      xlabel={Rendement brut du dividende sur douze mois (\%)},
      ylabel={Croissance annuelle du dividende sur \RegleAnneesCroissance~ans (\%)},
      x filter/.expression={x*1e2}, y filter/.expression={y*1e2},
      scatter, only marks, scatter src=explicit symbolic,
      scatter/classes={US={mark=*, blue}, DE={mark=square*, red}, FR={mark=triangle*, black}},
      legend pos=north east, legend cell align=left, legend style={fill=white}]
      % Classes dans l'ordre de la legende; un pays absent de scatter/classes ferait
      % echouer la compilation (portefeuille actuel: US, DE, FR seulement).
      \addplot table[x=rendite, y=wachstum, meta=land]{data/fig_rendite_wachstum.csv};
      \legend{\'Etats-Unis, Allemagne, France}
    \end{axis}
  \end{tikzpicture}
  \caption{Rendement brut du dividende et croissance annuelle moyenne du dividende sur
  \RegleAnneesCroissance~ans (nominale, en monnaie de cotation) pour chacun des \NombreTitres{}
  titres du portefeuille-exemple, par pays.}\label{fig:rendement-contre-croissance}
\end{figure}

Ce graphique montre que les titres \`a rendement \'elev\'e ont en g\'en\'eral cru moins vite~: le
coefficient de corr\'elation entre les deux grandeurs vaut \CorrelationRendementCroissance, une
relation n\'egative mais faible. La croissance va de \CroissanceTitreMin~\% par an
(\CroissanceTitreMinTicker) \`a \CroissanceTitreMax~\% (\CroissanceTitreMaxTicker), avec une m\'ediane
de \CroissanceMediane~\%. Les titres europ\'eens ont un rendement m\'edian de
\RenditeMedianeEurope~\% contre \RenditeMedianeUs~\% pour les am\'ericains, pour une croissance
m\'ediane voisine (\CroissanceMedianeEurope{} contre \CroissanceMedianeUs~\%).

\bk{Il serait tentant de projeter cette croissance pass\'ee dans l'avenir~; le mod\`ele s'en garde.} Ces
taux sont nominaux et mesur\'es sur des titres choisis justement parce que leur dividende a bien
progress\'e~: les prolonger reviendrait \`a supposer que les gagnants d'hier restent les gagnants de
demain. Le plan retient pour l'ETF maison une croissance r\'eelle de \CroissanceModeleMaison~\% par
an, celle qui laisse son rendement total au niveau du march\'e. La presse sp\'ecialis\'ee rappelle \`a
juste titre que des titres de qualit\'e portent la croissance d'un portefeuille \`a long
terme\QP{boerseonline_fire_portfolio}~; cela ne dispense pas de ce r\'ealisme.

\subsection{Deux dividendes r\'eels}\label{sec:porte-exemples}

La figure~\ref{fig:deux-dividendes-reels} suit ann\'ee apr\`es ann\'ee deux dividendes r\'eels~: celui
du titre retenu dont le dividende a mont\'e le plus longtemps sans interruption, et celui du titre de
tout l'univers qui a subi la plus forte baisse.

\annahme{Dividendes annuels par action, en monnaie de cotation, ramen\'es aux prix de 2026 avec
l'inflation allemande publi\'ee de 2019 \`a 2025\QH{inflation_destatis_2019_2025} et
\InflationZiel~\% par an hors de cette p\'eriode\QH{inflation_ziel}. Seules les ann\'ees compl\`etes
des donn\'ees sont prises en compte.}

\linienfigur{fig_div_beispiele}[xlabel={Ann\'ee}, ylabel={Dividende par action (prix de 2026)},
  xticklabel style={/pgf/number format/set thousands separator={}}, ymin=0,
  every axis plot post/.append style={mark=*, mark size=1.5pt},
  legend pos=north west, legend cell align=left, legend style={fill=white}]%
  {jahr}{titre_hausse,titre_baisse}%
  {{\BeispielHausseTicker},{\BeispielBaisse}}%
  {Dividende annuel r\'eel par action de \BeispielHausseTicker{}, titre retenu \`a la plus longue
  hausse continue, et de \BeispielBaisse, titre de l'univers \`a la plus forte baisse, en monnaie de
  cotation aux prix de 2026.\label{fig:deux-dividendes-reels}}

Ce graphique montre deux trajectoires oppos\'ees. Le dividende de \BeispielHausse{}
(\BeispielHausseTicker) a augment\'e chaque ann\'ee de \BeispielHausseDebut{} \`a \BeispielHausseFin,
soit \BeispielHausseHausses{} hausses d'affil\'ee et \BeispielHausseReelTotal~\% de pouvoir d'achat en
plus. M\^eme ce dividende exemplaire a pein\'e face \`a l'inflation~: en \BeispielHausseReelMinAnnee,
sa hausse de \BeispielHausseNominalMinAnnee~\% a \`a peine couvert une inflation de
\BeispielHausseInflationMinAnnee~\%, et sa progression r\'eelle n'a \'et\'e que de
\BeispielHausseReelMin~\%. La fen\^etre des donn\'ees commence en \DonneesDebut~; la s\'erie de hausses
peut \^etre plus longue, ce que ce document n'affirme pas faute de donn\'ees.

Le dividende de \BeispielBaisse, valeur du \BeispielBaisseIndice, a \'et\'e supprim\'e en
\BeispielBaisseAnnee{} (baisse de \BeispielBaisseProzent~\%), puis coup\'e de nouveau en
\BeispielBaisseAutreAnnee{} (\BeispielBaisseAutreProzent~\%). Sur les \BeispielBaisseAnneesSansBaisse{}
ann\'ees de donn\'ees qui pr\'ec\`edent la suppression, il n'avait jamais baiss\'e de plus de
\RegleBaisseSeuil~\%~: le crit\`ere \enquote{aucune baisse} ne l'aurait pas \'ecart\'e la veille de la
coupure. En \BeispielBaisseDerniereAnnee, il reste \BeispielBaisseEcartFinProzent~\%
\BeispielBaisseEcartFinSens{} son niveau r\'eel de \BeispielBaisseAnneeAvant. Ce titre
\BeispielBaisseStatut{} du portefeuille-exemple~: la r\`egle l'\'ecarte aujourd'hui, parce qu'elle
regarde en arri\`ere et voit la baisse. \bk{Aucune r\`egle fond\'ee sur le pass\'e ne prot\`ege d'une
coupure future~; seule la diversification en limite l'effet.} Le chapitre~\ref{sec:risques} montre que de telles
baisses ont frapp\'e l'ensemble du march\'e pendant les crises.

\subsection{Les biais de ce portefeuille}\label{sec:porte-biais}

Ce portefeuille-exemple porte quatre biais qu'il faut conna\^itre avant d'en tirer quoi que ce
soit.

Le premier est le biais du survivant. L'univers est compos\'e des membres actuels des indices,
c'est-\`a-dire des entreprises qui ont surv\'ecu et prosp\'er\'e. La liste des aristocrates du dividende
ne retient par d\'efinition que des soci\'et\'es qui augmentent leur dividende depuis des d\'ecennies~;
celles qui l'ont coup\'e en sont sorties et n'apparaissent pas dans l'univers. Leur historique est
donc flatteur par construction, et le titre en hausse continue de la
figure~\ref{fig:deux-dividendes-reels} en est un produit direct.

Le deuxi\`eme est l'emploi de donn\'ees courantes pour juger le pass\'e. Le rendement, le taux de
distribution, le flux de tr\'esorerie libre et le secteur sont ceux du \DateDonneesPortefeuille, et la
croissance est mesur\'ee sur les \RegleAnneesCroissance{} ann\'ees qui pr\'ec\`edent. Un investisseur
qui aurait appliqu\'e les m\^emes r\`egles en \DonneesDebut{} aurait choisi d'autres titres, et ce
document ne mesure pas comment ceux-l\`a se seraient comport\'es~: il n'existe aucun test historique
de la m\'ethode, et rien ne montre qu'elle bat un ETF monde.

Le troisi\`eme est l'univers limit\'e. Les \'Etats-Unis n'y sont repr\'esent\'es que par leurs
aristocrates du dividende, l'Europe par les grandes valeurs du DAX et de l'EURO STOXX~50~; ni le
Royaume-Uni, ni la Suisse, ni le Japon, ni les pays \'emergents, ni les moyennes capitalisations n'y
figurent. Le portefeuille ne compte finalement que \PortefeuillePays{} pays.

Le quatri\`eme est la courte p\'eriode des donn\'ees. \RegleHistoriqueAns{} ann\'ees de cours, de
\DonneesDebut{} \`a \DonneesFin, ne contiennent pas la crise financi\`ere de 2008~; la mesure du risque de la figure~\ref{fig:risque-nombre-titres} et
les croissances du tableau~\ref{tab:titres-portefeuille-exemple} en h\'eritent.

\bk{Retiens de ce chapitre qu'un ETF maison de \NombreTitres{} titres choisis par r\`egles
diversifie l'essentiel du risque propre \`a chaque titre, mais que la part que tu lui confies ne
d\'ecide pas de la date de ta retraite~: \`a \SparquoteZwanzig~\% d'\'epargne, aucune r\'epartition
n'atteint l'objectif, et la r\`egle des dividendes seuls reste en retrait du retrait de
\TauxRetrait~\% sur un ETF monde.} La m\'ethode de s\'election se r\'ep\`ete chaque ann\'ee avec des
donn\'ees fra\^iches~; les \NombreTitres{} titres de ce chapitre n'en sont qu'une photographie. Le
chapitre~\ref{sec:risques} met cette strat\'egie \`a l'\'epreuve de l'histoire~: baisses de dividendes pendant les
crises, inflation et ordre des rendements.

\Yannbox{\AnneeYannBasis}{Yann a vers\'e chaque mois la part maison de son \'epargne, selon la
r\'epartition \ReferenzStrategie, \`a parts \'egales dans les \NombreTitres{} titres de son ETF maison,
soit \PoidsParTitre~\% par titre~; ses titres sont pour \PartUs~\% am\'ericains, \PartDe~\% allemands et
\PartFr~\% fran\c{c}ais. Ce nombre de titres ram\`ene la fluctuation annuelle de son ETF maison de
\RisqueUnTitre~\% pour un titre seul \`a \RisqueTousTitres~\%. Si l'un de ses titres supprimait son
dividende comme \BeispielBaisse{} en \BeispielBaisseAnnee, ses dividendes maison baisseraient
d'environ \PoidsParTitre~\%, \`a rendement \'egal entre titres. S'il n'avait \'epargn\'e que
\SparquoteZwanzig~\% de son salaire, tout placer en maison lui aurait vers\'e
\RevenuAufteilungMaisonHundert~€ nets par mois en \AnneeHorizonFin, contre \RevenuHorizonZwanzig~€
avec la r\'epartition \ReferenzStrategie~: deux montants sous sa cible de \ZielMonat~€.}
